\section{Limitations}
\label{sec:limitations}
% Jujur & spesifik — reviewer menghargai ini. Sumber: caveat di findings.

\textbf{Scope of replication.} The representational findings replicate on
three checkpoints of a second model family (Sec.~\ref{sec:crossmodel}),
but the patching experiments---and the map-use dissociation with
them---have so far only run on Mistral-7B; output-side near-invariance
does replicate on an instruction-tuned Mistral (Sec.~\ref{sec:ceiling}),
so it is not a base-model artifact. Every patch is single-layer, so
redundant encoding across layers remains an alternative explanation for a
null result only partially ruled out (Sec.~\ref{sec:sweep}, caveat v).
\starhead{} itself was first noticed in per-type top-10 lists before
being fixed for the all-type test; Sec.~\ref{sec:crossmodel} supplies a
second-model check (L33\,H9, selection-free on five types), but a
preregistered replication is still outstanding.

\textbf{Measurement interface and contexts.} Output distributions are
read from option-letter logits after \texttt{Answer:}---the standard
silicon-sampling interface, but known to diverge from free-text answers
\citep{myanswerisc2024}. Our output-side claims (Sec.~\ref{sec:gap})
therefore scope to this interface: identity barely reaches the
letter-probability read-out, not necessarily free-text generation, which
we did not measure; the representation-side claims
(Sec.~\ref{sec:act2}) do not depend on the read-out at all. Relatedly,
the two sides are measured in different input contexts (identity-only
vs.\ QA), and the map is measurably weaker in the QA context
(Sec.~\ref{sec:sweep})---the fidelity--causality correlation is unchanged
there, but our headline fidelity numbers describe an identity-only prompt
and should be read as such.

\textbf{Ground truth and sample.} All distances are measured against one
US-centric institution, Pew ATP; cross-institution transfer (e.g.,
scoring the map against GSS-derived inter-group distances) and
cross-country replication (e.g., WVS, with a known language confound) are
future work. The dissociation claim rests on six attribute types---a
purposive sample of the 66 possible pairs among our twelve attributes
(Sec.~\ref{sec:setup}, Appendix~\ref{app:data})---enough for two
unambiguous counterexamples but not to generalize to intersections
outside this set, and some cells have thin survey samples, with NaN
pairs excluded (Appendix~\ref{app:data}). With cluster-robust inference
the causal experiments have twelve effective units per type: enough for
one clear result and direct between-type contrasts, not to certify a
sweep winner against 31 rivals or estimate the fidelity--causality
relationship precisely ($n{=}6$ types); scaling pairs, not questions, is
the binding constraint for follow-up work.

\textbf{Correlational map, one lexical encoder.} RSA fidelity is
correlational; the patching program (Sec.~\ref{sec:v2},
Sec.~\ref{sec:sweep}) addresses use, not the map's origin. The
partial-fidelity control (Sec.~\ref{sec:lexical}) uses one small sentence
encoder as the lexical baseline; static word vectors cannot represent
several value strings (e.g.\ age ranges), so a stronger lexical model
would only make the control stricter.

\section*{Ethics Statement}
Simulating the opinions of demographic groups is dual-use: it supports
legitimate survey-methodology research---pilot design, non-response
modelling, coverage checks---but the same capability can manufacture
apparently representative public opinion, or reinforce stereotypes by
presenting a model's guess about a group as that group's view.

Our findings cut against the permissive reading: current LLM outputs are
\emph{not} substitutes for human respondents. Under natural prompting the
model's answers move by less than 2\% of their error when the entire
demographic identity is replaced, and the small movement that exists is
not reliably toward the target group's truth. The failure is also
unequal---race-based attributes are the weakest and most prompt-fragile
everywhere measured, so simulated outputs for racial groups are least
trustworthy where misuse would be most harmful. We therefore caution
against policy-adjacent uses of simulated group opinions.

Interpretability results of this kind carry their own risk: a faithful
location is a natural target for steering, and Sec.~\ref{sec:v2} shows
naive steering through such a location can push predictions \emph{away}
from the truth in low-fidelity types---read the map before trusting it as
a lever.

All data are public, aggregated Pew ATP response distributions with no
personally identifiable information; no human subjects were recruited for
this work.
